\documentclass[12pt]{article}

\usepackage{sbc-template}
\usepackage{graphicx,url}
\usepackage[utf8]{inputenc}
\usepackage[T1]{fontenc}
\usepackage[brazil]{babel}
\usepackage{booktabs,tabularx,array}
\usepackage{tikz}
\usetikzlibrary{arrows.meta,positioning,fit,backgrounds}
\usepackage{pgfplots}
\pgfplotsset{compat=1.18}

\sloppy
\setlength{\tabcolsep}{4pt}
\newcommand{\cod}[1]{\texttt{\small #1}}

\title{Plataforma Distribuída de Busca e Indexação\\ do Acervo de Provas do ENADE}

\author{Gustavo Barros Sampaio\inst{1}}
\address{Curso de Sistemas de Informação -- Centro Universitário Católica de Quixadá\\
  Disciplina de Sistemas Distribuídos}

\begin{document}

\maketitle

\begin{abstract}
This paper describes the design, implementation and evaluation of a distributed platform that makes the
ENADE exam archive searchable. Official PDFs published by INEP (84 exams, 2021--2023) are parsed into 3,360
question documents with source page, answer key and metadata; documents flow through Apache Kafka to
Elasticsearch, and a FastAPI service with Redis cache serves a web client. Users authenticate against OpenLDAP
and services authenticate mutually with Kerberos. We measured p95 search latency of 19.8\,ms without cache
and 6.6\,ms with cache, and no message loss or duplication after induced crashes.
\end{abstract}

\begin{resumo}
Este artigo descreve o projeto, a implementação e a avaliação de uma plataforma distribuída que torna
pesquisável o acervo do ENADE. PDFs oficiais do INEP (84 provas, 2021--2023) são convertidos em 3.360
documentos de questão com página de origem, gabarito e metadados; os documentos trafegam pelo Apache Kafka até
o Elasticsearch, e uma API FastAPI com cache Redis atende um cliente web. Usuários autenticam-se no OpenLDAP e
os serviços, mutuamente, via Kerberos. Medimos latência p95 de busca de 19,8\,ms sem cache e 6,6\,ms com
cache, e nenhuma perda ou duplicação após falhas induzidas.
\end{resumo}

\section{Introdução}

O ENADE acumula, desde 2004, provas, gabaritos e padrões de resposta de dezenas de cursos, publicados em PDFs
heterogêneos \cite{inepprovas}. O objetivo deste trabalho é transformar esse acervo em uma base pesquisável,
resiliente e segura, seguindo a arquitetura de cinco camadas do enunciado (ingestão, Kafka, Elasticsearch,
API/Redis e cliente web) e uma camada transversal de segurança (Kerberos entre serviços e LDAP para usuários).
O trabalho é individual; todo o código, configurações e evidências estão no repositório
(\cod{docker-compose.yml}, \cod{codigo/}, \cod{evidencias/}). As métricas deste texto foram produzidas por
execução real no ambiente descrito na Seção~\ref{sec:avaliacao}; requisitos não atendidos integralmente são
apontados explicitamente.

\section{Arquitetura}\label{sec:arq}

A Figura~\ref{fig:arq} mostra os serviços (todos em Docker Compose, rede \cod{enade-net}, domínio interno
\cod{enade.local}). O \textbf{ingestor} coleta e extrai os PDFs e publica um JSON por questão no tópico
\cod{enade.questoes.v1}; o \textbf{indexer} consome, valida e indexa em lote no Elasticsearch, desviando
mensagens inválidas para \cod{enade.questoes.dlq}; a \textbf{API} consulta o Elasticsearch, usa o Redis para
cache, sessão e carrinho, autentica usuários no OpenLDAP e serve a interface web.

\begin{figure}[ht]
\centering
\resizebox{\textwidth}{!}{%
\begin{tikzpicture}[font=\sffamily\footnotesize,
  caixa/.style={draw, rounded corners=2pt, align=center, text width=25mm, minimum height=12mm, fill=white, inner sep=2pt},
  seta/.style={-{Latex[length=1.8mm]}, thick}, rot/.style={font=\sffamily\scriptsize, fill=white, inner sep=1pt}]
\node[caixa] (inep) at (0,0) {PDFs oficiais\\do INEP};
\node[caixa] (ing) at (3.25,0) {ingestor\\(coleta, extração,\\produtor)};
\node[caixa] (kafka) at (6.5,0) {Kafka 3.9 (KRaft)\\questoes.v1: 3 part.\\dlq: 1 part.};
\node[caixa] (idx) at (9.75,0) {indexer\\(consumidor,\\DLQ, commit)};
\node[caixa] (es) at (13,0) {Elasticsearch\\8.19, índice\\PT-BR};
\node[caixa] (redis) at (9.75,-2.4) {Redis\\(cache, sessão,\\carrinho)};
\node[caixa] (api) at (13,-2.4) {API FastAPI\\+ interface web};
\node[caixa] (ldap) at (9.75,-4.6) {OpenLDAP\\(usuários,\\grupos)};
\node[caixa] (nav) at (13,-4.6) {Navegador\\(Estudante,\\Professor)};
\node[caixa, text width=52mm] (kdc) at (3.25,-2.6) {MIT Kerberos KDC\\realm ENADE.LOCAL\\(um keytab por serviço)};
\draw[seta] (inep) -- (ing);
\draw[seta] (ing) -- (kafka);
\draw[seta] (kafka) -- (idx);
\draw[seta] (idx) -- (es);
\draw[seta] (api) -- node[rot, right=1pt]{SPNEGO} (es);
\draw[seta, <->] (api) -- (redis);
\draw[seta] (idx) -- node[rot, left=1pt]{versão} (redis);
\draw[seta] (api) -- node[rot, sloped, above=1pt]{bind} (ldap);
\draw[seta, <->] (nav) -- node[rot, right=1pt]{HTTP} (api);
\draw[seta, dashed, gray] (kdc) -- node[rot, left=1pt]{tickets} (ing);
\draw[seta, dashed, gray] (kdc.north east) -- (kafka.south);
\end{tikzpicture}}
\caption{Arquitetura implantada. As setas do pipeline são autenticadas por Kerberos: SASL/GSSAPI entre ingestor, Kafka e indexer; SPNEGO (com autenticação mútua) do indexer e da API para o Elasticsearch. Tracejado: emissão de tickets pelo KDC (todos os serviços usam seus keytabs).}
\label{fig:arq}
\end{figure}

\textbf{Topologia e justificativas.} O laboratório dispõe de 4 vCPUs e 5,8\,GB de RAM para o Docker, o que
limita o cluster a um nó de cada componente (decisão ADR-01). O tópico principal tem \textbf{3 partições}: a
chave da mensagem é o \textit{id} determinístico da questão, o que preserva a ordem das versões de uma mesma
questão e permite até três consumidores paralelos no grupo \cod{enade-indexer}; a DLQ tem 1 partição (baixo
volume, ordem de chegada útil para diagnóstico). O índice \cod{enade-questoes-v1} tem \textbf{1 shard e 0
réplicas}: o acervo ocupa poucos MB, muito abaixo da faixa usual de 10--50\,GB por shard, e réplicas não podem
ser alocadas com um único nó; o alias \cod{enade-questoes} permite reindexar para um índice com mais shards
sem alterar a API. Em produção, a mesma configuração escalaria horizontalmente com 3 brokers
(\cod{replication.factor=3}, \cod{min.insync.replicas=2}), 3 nós Elasticsearch com 1 réplica e réplicas da
API sem estado (sessão e cache já estão no Redis).

\section{Coleta e Engenharia de Dados}\label{sec:coleta}

\textbf{Coleta (RF01).} O coletor lê as páginas anuais do INEP \cite{inepprovas}, descobre os links de prova
(\cod{PV}), gabarito (\cod{GB}) e padrão de resposta, descarta versões acessíveis duplicadas (ampliada, ledor)
e registra URL, tamanho, SHA-256 e data. Dois problemas reais foram tratados sem atalhos: o servidor de
download envia cadeia TLS incompleta (adicionamos a CA intermediária pública da RNP ao \textit{bundle}, mantendo
a verificação) e, em 2022, os gabaritos têm nomes diferentes das provas (pareamento pela ordem da página). Foram
coletados 252 PDFs (972\,MB) de 2021, 2022 e 2023; as demais edições não foram processadas por limite de tempo e
por usarem layouts distintos.

\textbf{Extração (RF02--RF04, RF06).} A extração é determinística, com PyMuPDF. As linhas são reconstruídas a
partir das posições de cada caractere, o que permitiu corrigir: (i) espaços espúrios gerados pelas ligaduras
``ti''/``fi'' da fonte Calibri (``Insti tuto''), identificados por se sobreporem ao glifo anterior; (ii) a
codificação da fonte Symbol (``Î''$\to\in$, ``Û''$\to\Leftrightarrow$); (iii) índices e expoentes, detectados por
tamanho de fonte e deslocamento da linha de base e convertidos em LaTeX inline (ex.: \cod{\$Z\_\{k-1\}\$});
(iv) páginas em duas colunas (2021/2022). Cabeçalhos \textit{QUESTÃO NN} e \textit{QUESTÃO DISCURSIVA NN}
delimitam as questões; o componente (FG/CE) vem da tabela de distribuição da capa ou do título de seção; as
alternativas são encontradas por busca reversa da sequência A--E (ou A--D, formato real de Medicina 2023);
o texto-base é separado do enunciado pela última referência bibliográfica. Tabelas com borda viram Markdown e
figuras (imagens e agrupamentos de desenhos vetoriais) são recortadas como PNG. Cada documento guarda o PDF, a
URL e as páginas de origem.

\textbf{Metadados (RF05).} Área (mapeamento próprio do curso para a grande área), curso (lido do cabeçalho
corrente da prova), ano, grau e componente. Os PDFs não informam as competências de cada questão; usamos a
taxonomia oficial da Portaria Inep 289/2023 (8 competências e 13 temas de Formação Geral)
\cite{inep2023portaria289} e um classificador léxico explícito, marcado no documento como \textit{inferido}.
A mesma taxonomia é aplicada às três edições para permitir séries temporais comparáveis.

\textbf{LLM.} Nenhum LLM foi usado no pipeline de dados (custo zero e resultado reprodutível). Consequência:
diagramas não são convertidos para Mermaid; ficam como imagem (RF04 parcial).

\begin{table}[ht]
\centering
\caption{Volume do acervo extraído (execução de 22/11/2026).}\label{tab:volume}
\small
\begin{tabular}{lrrrr}
\toprule
 & 2021 & 2022 & 2023 & Total \\
\midrule
Provas (cursos) & 30 & 26 & 28 & 84 \\
PDFs coletados (prova+gabarito+padrão) & 90 & 78 & 84 & 252 \\
Questões & 1.200 & 1.040 & 1.120 & 3.360 \\
\quad objetivas / discursivas & 1.050 / 150 & 910 / 130 & 1.064 / 56 & 3.024 / 336 \\
\quad anuladas (sinalizadas) & 7 & 2 & 5 & 14 \\
\quad com figura recortada & 375 & 409 & 408 & 1.192 \\
Grupos de FG idênticos entre cursos & 10 & 10 & 10 & 30 \\
\bottomrule
\end{tabular}
\end{table}

A Tabela~\ref{tab:volume} resume o acervo. Todas as 3.024 objetivas receberam gabarito ou marca de anulação a
partir do gabarito oficial; 281 questões têm fórmulas em LaTeX e 280 têm tabelas em Markdown. A última linha
mostra que as 10 questões de Formação Geral de cada ano, repetidas em todos os cursos daquele ano, foram
reconhecidas como idênticas pelo \textit{hash} de conteúdo --- base da busca interdisciplinar (Seção~\ref{sec:busca}).

\section{Ingestão Kafka/Elasticsearch}\label{sec:ingestao}

\textbf{Produtor (RF07).} Publica cada questão com chave = \textit{id}, \cod{acks=all},
\cod{enable.idempotence=true} e compressão zstd. O mesmo esquema Pydantic (com \cod{extra=forbid}) valida no
produtor e no consumidor.

\textbf{Consumidor e resiliência (RF08, RF09).} O indexer usa \cod{enable.auto.commit=false}: consome lotes de
até 200 mensagens, envia um \textit{bulk} ao Elasticsearch e só então confirma o offset (síncrono). Falhas
transitórias (conexão, 429, 5xx) disparam novas tentativas com \textit{backoff} exponencial, sem confirmar o
offset. Assim, a entrega é \textit{at-least-once}; somada à idempotência, o efeito no índice é de
\textit{exactly-once}.

\textbf{Idempotência (RF10).} O \textit{id} é o SHA-256 de (ano, curso, tipo, número), usado como chave Kafka e
\cod{\_id} no Elasticsearch: reprocessar sobrescreve o documento em vez de duplicá-lo.

\textbf{DLQ (RF11).} JSON inválido, esquema inválido (campo ausente, campo desconhecido, \textit{id}
inconsistente) ou rejeição 4xx do mapping vão para a DLQ com cabeçalhos de diagnóstico (motivo, tópico,
partição e offset de origem); o offset só é confirmado depois que a DLQ confirma a escrita.

\textbf{Modelagem Elasticsearch.} Mapping explícito com \cod{dynamic: strict}. Campos textuais (texto-base,
enunciado, alternativas) usam o analisador de indexação \cod{lowercase} $\to$ \cod{stop}
(\cod{\_brazilian\_}) $\to$ \cod{asciifolding} $\to$ \cod{stemmer light\_portuguese} \cite{savoy2006light}; o
analisador de busca acrescenta \cod{synonym\_graph} (dicionário próprio, recarregável) \emph{depois} do
\textit{stemmer} --- ordem corrigida após um teste mostrar que ``desigualdade'' e ``desigualdades'' expandiam
sinônimos de forma diferente. Os campos têm \cod{term\_vector} com offsets para \textit{highlight} rápido e
\cod{copy\_to} para \cod{texto\_completo} (usado em \textit{More Like This}). Filtros e agregações usam campos
\cod{keyword}; o gabarito e o padrão de resposta não são indexados para busca textual, evitando que um
Estudante infira a resposta por meio de consultas.

\section{Busca e API}\label{sec:busca}

A API (FastAPI, documentação OpenAPI em \cod{/docs}) expõe, entre outras: \cod{GET /api/busca} (texto, filtros
combináveis por área, curso, ano, componente, tipo, competência, tema e grau; paginação; \textit{highlight};
facetas), \cod{GET /api/questoes/\{id\}}, \cod{/gabarito} (Professor), \cod{/equivalentes},
\cod{/api/carrinho} (lista), \cod{/api/carrinho/exportar} (PDF/JSON, Professor), \cod{/api/analitico}
(Professor), \cod{/health/ready} e \cod{/metrics}. Os códigos HTTP seguem a semântica: 201 ao incluir na lista,
204 em remoções, 400/422 para parâmetros inválidos, 401 sem sessão, 403 sem permissão, 404, 409 (lista cheia),
429 (bloqueio após 5 falhas de login) e 503 (dependência indisponível).

\textbf{Full-text e relevância (RF13).} \cod{multi\_match} \textit{best\_fields} com pesos (enunciado$^3$,
texto-base$^{1{,}5}$, alternativas) e reforço por frase (\textit{phrase}, \textit{slop} 2). \textbf{Interdisciplinaridade
(RF14).} Uma mesma busca percorre todos os cursos; a opção ``agrupar'' usa \textit{field collapsing} por grupo de
equivalência, e a rota \cod{/equivalentes} retorna a mesma questão em outros cursos e as mais semelhantes de
outros cursos (\textit{More Like This}). Por exemplo, ``ética'' retorna 179 questões de 67 cursos diferentes e
destaca também ``moral'' (sinônimo); ``idadismo'', que não ocorre no acervo, recupera 28 questões sobre
etarismo via dicionário de sinônimos. \textbf{Agregações (RF15).} Histograma de competências e de temas por ano
(\textit{terms} + \textit{histogram} sobre o ano) e distribuição de questões por curso (com FG/CE e anuladas),
exibidos no painel analítico.

\section{Cache e Sessão}\label{sec:cache}

\textbf{Cache (RF18).} A chave é \cod{cache:busca:v\{versão\}:\{perfil\}:\{sha256(parâmetros)\}} com TTL de 300\,s.
A invalidação é por versão: o indexer incrementa \cod{indice:versao} no Redis a cada lote indexado, de modo que
respostas antigas nunca são reutilizadas após uma atualização e expiram pelo TTL; Professores também podem
invalidar manualmente. Incluir o perfil na chave impede que uma resposta com gabarito seja servida a um
Estudante (há teste para isso). Se o Redis cair, a busca degrada para consulta direta. \textbf{Sessão e
carrinho (RF19).} Após o login, um token aleatório de 256 bits é gravado no Redis com TTL deslizante de 1\,h e
enviado em \textit{cookie} \cod{HttpOnly}/\cod{SameSite=Strict}; o carrinho é um \textit{sorted set} por sessão
(ordem de inclusão, limite de 200). \textbf{Exportação (RF21).} PDF (ReportLab) com texto, figuras recortadas,
origem e folha de respostas, ou JSON completo.

\section{Segurança}\label{sec:seguranca}

\textbf{Kerberos (RF22).} Um KDC MIT (realm \cod{ENADE.LOCAL}) é criado na subida e exporta um
\textit{keytab} por serviço, montado apenas no contêiner correspondente (Tabela~\ref{tab:principals}). No
Kafka, todos os listeners (clientes, inter-broker e controller KRaft) usam SASL/GSSAPI \cite{kafkadocs}, com
\cod{StandardAuthorizer} e ACLs mínimas. No Elasticsearch, o realm \cod{kerberos} valida tokens SPNEGO
\cite{elastickerberos} e \textit{role mappings} associam cada principal a um papel. O cliente exige
autenticação mútua (\cod{mutual\_authentication=REQUIRED}): o servidor precisa provar sua identidade na
resposta \cite{neuman1994kerberos}.

\begin{table}[ht]
\centering
\caption{Principals Kerberos e permissões efetivas.}\label{tab:principals}
\small
\begin{tabularx}{\textwidth}{lX}
\toprule
Principal & Uso e permissão \\
\midrule
\cod{kafka/kafka.enade.local} & identidade do broker/controller (super usuário do cluster) \\
\cod{HTTP/elasticsearch.enade.local} & aceitar SPNEGO no Elasticsearch \\
\cod{ingestor} & Kafka: \textit{Write} em questoes.v1 e DLQ \\
\cod{indexer} & Kafka: \textit{Read} questoes.v1 e grupo; \textit{Write} DLQ. ES: papel \cod{enade\_indexer} (escrita) \\
\cod{api} & ES: papel \cod{enade\_api} (somente leitura). Kafka: \textit{Describe} (lag) e leitura da DLQ \\
\cod{admin} & criação de tópicos e ACLs (job de setup) \\
\cod{intruso} & principal válido sem permissões (testes negativos) \\
\bottomrule
\end{tabularx}
\end{table}

\textbf{LDAP e perfis (RF23--RF25).} A árvore \cod{dc=enade,dc=local} tem \cod{ou=usuarios} e \cod{ou=grupos}
(\cod{cn=estudantes}, \cod{cn=professores}), com usuários sintéticos de teste. O login faz \textit{bind} direto
com o DN do usuário (a API não armazena senhas) e o perfil vem do grupo, consultado com uma conta somente
leitura; o \textit{uid} é validado para impedir injeção LDAP. A autorização é feita no backend: para Estudante,
\cod{gabarito} e \cod{padrao\_resposta} são excluídos de todas as respostas (busca, detalhe, carrinho) e as rotas
de gabarito, painel analítico, exportação e operação retornam \textbf{HTTP 403}.

\textbf{Testes.} Sete testes automatizados de Kerberos passam: emissão de TGT (\cod{klist}); Elasticsearch
rejeita requisição sem credencial (401), aceita SPNEGO mútuo do \cod{indexer} e nega (403) o principal
\cod{intruso}, que autentica mas não tem papel; o Kafka não fornece metadados a cliente sem SASL e recusa a
publicação do \cod{intruso} (\cod{TOPIC\_AUTHORIZATION\_FAILED}). Trinta testes de integração da API cobrem 401,
injeção LDAP, ausência de gabarito para Estudante, 403 em sete rotas restritas e a não contaminação do cache
entre perfis. \textbf{Segredos (RNF08):} \cod{.env} e keytabs ficam fora do Git, cada contêiner recebe apenas
as variáveis de que precisa e cada \textit{commit} passa por varredura das senhas do \cod{.env}.

\section{Avaliação Experimental}\label{sec:avaliacao}

\textbf{Ambiente.} Windows 11 + Docker Desktop (WSL2), 4 vCPUs e 5,8\,GB de RAM para os contêineres; um nó de
cada serviço; acervo de 3.360 documentos. Os números abaixo vêm dos arquivos em \cod{evidencias/}.

\textbf{Qualidade da extração (RNF03).} Indicadores estruturais automáticos: 40 questões por prova em todas as
84 provas, 0 documentos inválidos, 3.318 questões (98,8\%) sem alertas, 29 com quatro alternativas
(Medicina 2023, formato real) e 11 (0,3\%) sem alternativas detectadas. Esses indicadores \emph{não}
substituem a validação humana exigida: foi sorteada uma amostra estratificada e reprodutível de 60 questões
(semente 2026) e construída a tela \cod{/validacao}, que mostra a página original do PDF ao lado do conteúdo
extraído com seis critérios; \textbf{a taxa de acerto manual ainda não foi medida} e será incluída quando a
revisão humana for concluída.

\textbf{Vazão de ingestão e lag (RNF02).} A extração levou 641--728\,s para as 84 provas (7,6--8,7\,s por
prova, um núcleo). O produtor publicou as 3.360 mensagens em 3,2\,s (1.046\,msg/s; 2.645\,msg/s com conexões
aquecidas). A Figura~\ref{fig:lag} mostra a carga inicial: o lag do grupo sobe a 3.360 quando a publicação
termina e cai a zero em cerca de 10\,s; o indexer processou 3.360 documentos em 17 lotes, $\approx$313\,docs/s de tempo
efetivo de processamento ($\approx$371\,docs/s de relógio entre o primeiro e o último lote). Do início da publicação ao
lag zero foram 14\,s (variando entre 14 e 25\,s em quatro execuções, conforme a carga da máquina). O
\textit{refresh} de 1\,s explica o atraso entre lag e contagem visível.

\begin{figure}[ht]
\centering
\begin{tikzpicture}
\begin{axis}[width=0.86\textwidth, height=5.2cm, xlabel={tempo desde o início da medição (s)},
  ylabel={mensagens / documentos}, legend style={font=\footnotesize, at={(0.5,1.03)}, anchor=south, legend columns=2},
  xmin=0, xmax=23, ymin=0, ymax=3700, grid=major, grid style={gray!25}, ticklabel style={font=\footnotesize}]
\addplot[thick, red!70!black, mark=*, mark size=1.2pt] coordinates
 {(0.0,0) (3.2,0) (4.7,0) (6.1,0) (7.6,3360) (9.1,3360) (10.4,3360) (11.8,3160) (13.2,3160) (14.6,2560)
  (16.0,1960) (17.4,760) (18.8,360) (20.1,0) (21.5,0) (22.8,0)};
\addlegendentry{lag do consumidor (mensagens)}
\addplot[thick, blue!70!black, dashed, mark=square*, mark size=1.2pt] coordinates
 {(0.0,0) (3.2,0) (4.7,0) (6.1,0) (7.6,0) (9.1,0) (10.4,1) (11.8,200) (13.2,200) (14.6,759) (16.0,1405)
  (17.4,2400) (18.8,2787) (20.1,3200) (21.5,3360) (22.8,3360)};
\addlegendentry{documentos visíveis no índice}
\end{axis}
\end{tikzpicture}
\caption{Carga inicial do acervo (amostragem a cada $\approx$1,4\,s via \cod{/api/admin/ingestao}).}
\label{fig:lag}
\end{figure}

\textbf{Idempotência e DLQ.} Republicar o mesmo acervo adicionou 3.360 mensagens ao tópico e variou em
\textbf{0} o número de documentos. Cinco mensagens malformadas injetadas (não-JSON, JSON não-objeto, campos
ausentes, \textit{id} inconsistente e campo desconhecido) foram para a DLQ com o motivo correto; o índice
permaneceu com 3.360 documentos e o consumo seguiu.

\textbf{Resiliência (RNF04).} O script \cod{experimento-falha.sh} recria o índice, reposiciona o grupo no início
do tópico (10.085 mensagens: três versões do acervo e as cinco inválidas) e: (1) mata o indexer com
\cod{SIGKILL} quando o índice tem 2.056 documentos --- nesse instante havia 4.600 offsets confirmados; (2)
reinicia o indexer e mata o Elasticsearch por 30\,s --- o lag ficou parado em 5.485 e o indexer fez 6
retentativas sem confirmar offsets; (3) religa o Elasticsearch. Em 144\,s o lag voltou a zero com
\textbf{3.360 de 3.360 documentos}: nenhuma perda e nenhuma duplicata. Efeito colateral esperado do
\textit{at-least-once}: no reprocessamento as cinco mensagens inválidas voltaram à DLQ (10 registros).

\textbf{Latência de busca e cache (RNF01).} Teste k6 com taxa constante de 20\,req/s por 60\,s em cada cenário,
180 consultas distintas (30 termos $\times$ 6 combinações de filtros), usuário Estudante. ``Sem cache'' força
\cod{cache=false}; ``com cache'' roda após um aquecimento. A Figura~\ref{fig:lat} mostra os percentis; ambos
ficam bem abaixo das metas (p95 $<$ 500\,ms e $<$ 50\,ms), com 0 erros em 2.583 requisições. O
\textbf{hit rate} foi 100\% nessa rodada; numa rodada anterior foi 85,3\% porque as 176 chaves aquecidas
expiraram (TTL) no meio da fase --- exatamente 176 \textit{misses} --- o que confirma a política de expiração.

\begin{figure}[ht]
\centering
\begin{tikzpicture}
\begin{axis}[ybar, width=0.8\textwidth, height=5cm, bar width=14pt, ymin=0, ymax=30,
  ylabel={latência (ms)}, symbolic x coords={p50,p95,p99}, xtick=data, enlarge x limits=0.25,
  legend style={font=\footnotesize, at={(0.5,1.03)}, anchor=south, legend columns=2}, nodes near coords,
  nodes near coords style={font=\scriptsize}, ticklabel style={font=\footnotesize}, grid=major, grid style={gray!20}]
\addplot[fill=red!55] coordinates {(p50,14.4) (p95,19.8) (p99,24.7)};
\addplot[fill=blue!50] coordinates {(p50,4.1) (p95,6.6) (p99,11.7)};
\legend{sem cache (Elasticsearch), com cache (Redis)}
\end{axis}
\end{tikzpicture}
\caption{Latência de \cod{GET /api/busca} medida pelo k6 (n = 1.201 requisições por cenário).}
\label{fig:lat}
\end{figure}

\section{Gargalos e Lições}\label{sec:licoes}

(1) \textbf{Biblioteca não thread-safe derrubou a API.} Na primeira rodada do k6, 349 requisições falharam e a
latência chegou a 43\,s: a API abortava com \textit{double free or corruption}. A causa era um único objeto de
autenticação SPNEGO (\cod{requests-gssapi}) compartilhado entre as threads do FastAPI. A correção --- um cliente
Elasticsearch por thread e envio oportunista do token --- eliminou as falhas (0 erros, p95 19,8\,ms). Lição:
segurança forte em todas as chamadas tem custo e precisa de testes de concorrência. (2) \textbf{Kerberos em
Python:} o \textit{wheel} do \cod{confluent-kafka} não traz GSSAPI; compilamos o librdkafka com Cyrus SASL. (3)
\textbf{Heterogeneidade dos PDFs:} duas colunas, quatro alternativas, nomes divergentes de arquivos e rótulos
ocultos de diagramação (um ``MATÉRIA'' de rodapé que vazava para a alternativa E) só apareceram ao processar o
acervo inteiro; os indicadores agregados (grupos de FG por ano, alertas por prova) foram decisivos para
encontrá-los. (4) \textbf{Ordem dos filtros} do analisador muda a semântica dos sinônimos. (5) O gargalo da
ingestão é a extração (um núcleo, $\approx$8\,s por prova), não o Kafka nem o Elasticsearch; ela poderia ser
paralelizada por arquivo. (6) \textbf{Limitações:} sem TLS interno (Kerberos autentica, mas não cifra), licença
\textit{trial} do realm Kerberos do Elasticsearch (30 dias), um nó por componente, diagramas sem Mermaid e
competências inferidas sem acurácia medida.

\section{Conclusão}

\begin{table}[ht]
\centering
\caption{Situação dos requisitos ao final do trabalho (evidências em \cod{docs/STATUS.md}).}\label{tab:status}
\small
\begin{tabularx}{\textwidth}{lX}
\toprule
Situação & Requisitos \\
\midrule
Concluído e testado & RF02, RF03, RF06--RF25; RNF01, RNF02, RNF04, RNF05, RNF06, RNF08 \\
Parcial & RF01 (só 2021--2023), RF04 (diagramas sem Mermaid), RF05 (competências inferidas, sem acurácia medida);
          RNF07 (roteiro não repetido em máquina limpa) \\
Pendente & RNF03 (validação manual da amostra de 60 questões) \\
\bottomrule
\end{tabularx}
\end{table}

A plataforma cumpre o fluxo completo PDF $\to$ extração $\to$ Kafka $\to$ Elasticsearch $\to$ API/Redis $\to$
web com dados reais e rastreáveis do INEP, autenticação mútua Kerberos entre serviços, login LDAP e perfis com
negação 403 verificada no backend. Medimos latência p95 de 19,8\,ms (sem cache) e 6,6\,ms (com cache), vazão de
centenas de documentos por segundo e ausência de perda/duplicação sob falhas induzidas. Ficam como trabalhos
futuros: concluir a validação manual da extração, ampliar o acervo para todas as edições, TLS interno,
conversão de diagramas assistida por LLM com validação e um cluster multinó.

\section{Divisão de Tarefas}

Trabalho individual: o(a) autor(a) foi responsável por todas as etapas (requisitos, arquitetura, implementação,
testes, medições e redação). \textbf{Uso declarado de IA generativa:} um assistente de programação baseado em
LLM (Claude Opus 5.5, Anthropic) foi usado como ferramenta de desenvolvimento para gerar e revisar código,
scripts e documentação, sob supervisão do autor; o custo desse uso não foi medido. Nenhum LLM foi usado no
pipeline de extração (custo zero), por isso não há taxa de erro de LLM a reportar; a qualidade da extração é
avaliada pelos indicadores da Seção~\ref{sec:avaliacao} e pela validação manual em andamento.

\bibliographystyle{sbc}
\bibliography{referencias}

\end{document}
